\section{Results}
\label{sec:results}

\subsection{Baseline versus full evidence chain}

Table~\ref{tab:main-comparison} reports the main baseline comparison. The full evidence chain improves every measured dimension over the trace-oriented baseline. Intent capture, policy checking, and receipt export all move from zero observed positive reviews in the baseline to 16 of 16 tasks in the evidence chain. The repository metric records \emph{execution verified} as 11 of 16, but that count should be read in two layers: among the 13 runs that reached the execution stage, 11 completed with verification success and 2 ended as explicit tamper-detected outcomes; the remaining 3 tasks are reviewable policy-blocked outcomes rather than missing bundles.


\begin{table}[t]
\centering
\caption{Trace-oriented baseline versus full evidence chain. Counts report positive bundle-review verdicts across the 16-task suite; scores are average rule-based metrics on a 0--5 scale.}
\label{tab:main-comparison}
\small
\resizebox{\linewidth}{!}{%
\begin{tabular}{lccccccccc}
\toprule
Mode & Intent & Policy & Exec.\ verified & Receipt & Explicitness & Replayability & Tamper & Audit & Integration \\
\midrule
    Baseline & 0/16 & 0/16 & 0/16 & 0/16 & 1.0 & 3.0 & 0.0 & 2.0 & 1.0 \\
    Evidence chain & 16/16 & 16/16 & 11/16 & 16/16 & 5.0 & 5.0 & 4.81 & 5.0 & 5.0 \\
\\
\bottomrule
\end{tabular}%
}
\end{table}


The metric changes are large. Average explicitness rises from 1.0 to 5.0, replayability from 3.0 to 5.0, tamper sensitivity from 0.0 to 4.81, audit boundedness from 2.0 to 5.0, and evidence-stage exposure from 1.0 to 5.0. This pattern supports the paper's narrower claim: within a stable exported-bundle contract, the evidence-chain surface preserves review questions that the trace-oriented surface does not explicitly retain.

\subsection{External baseline comparison}

Table~\ref{tab:external-comparison} compares the minimal live CrewAI baseline with the full evidence chain. The external baseline improves replayability relative to the local baseline, reaching an average score of 4.0, because a real CrewAI runtime now produces the task execution surface that is later normalized into the five-file bundle. However, it still leaves explicit intent capture, policy persistence, integrity verification, and receipt export absent. As a result, explicitness, tamper sensitivity, audit boundedness, and evidence-stage exposure remain at 1.0, 0.0, 2.0, and 1.0 respectively.


\begin{table}[t]
\centering
\caption{CrewAI-like external baseline versus the full evidence chain. The external baseline preserves a richer default execution surface than the local trace baseline, but does not add explicit evidence-chain semantics.}
\label{tab:external-comparison}
\small
\resizebox{\linewidth}{!}{%
\begin{tabular}{lccccccccc}
\toprule
Mode & Intent & Policy & Exec.\ verified & Receipt & Explicitness & Replayability & Tamper & Audit & Integration \\
\midrule
    External baseline & 0/16 & 0/16 & 0/16 & 0/16 & 1.0 & 4.0 & 0.0 & 2.0 & 1.0 \\
    Evidence chain & 16/16 & 16/16 & 11/16 & 16/16 & 5.0 & 5.0 & 4.81 & 5.0 & 5.0 \\
\\
\bottomrule
\end{tabular}%
}
\end{table}


This comparison only partially addresses the straw-baseline concern. Because the baseline is a minimal live integration rather than a full production benchmark, the result should be read as evidence that a framework-native default logging surface still leaves major evidence gaps by default, not as a definitive benchmark against every production framework.

\subsection{Ablation results}

Table~\ref{tab:ablation-comparison} summarizes the seven-mode comparison. The ablations show targeted degradation rather than uniform collapse.


\begin{table}[t]
\centering
\caption{Seven-mode comparison across the baseline, external baseline, four ablations, and the full evidence chain. Scores are average 0--5 values; counts report positive review verdicts across 16 tasks.}
\label{tab:ablation-comparison}
\small
\resizebox{\linewidth}{!}{%
\begin{tabular}{lccccccccc}
\toprule
Mode & Explicitness & Replayability & Tamper & Audit & Stage exposure & Intent & Policy & Exec.\ verified & Receipt \\
\midrule
    Baseline & 1.0 & 3.0 & 0.0 & 2.0 & 1.0 & 0/16 & 0/16 & 0/16 & 0/16 \\
    External baseline & 1.0 & 4.0 & 0.0 & 2.0 & 1.0 & 0/16 & 0/16 & 0/16 & 0/16 \\
    No intent & 3.0 & 4.0 & 4.81 & 4.0 & 4.0 & 0/16 & 16/16 & 11/16 & 16/16 \\
    No governance & 2.0 & 5.0 & 5.0 & 4.0 & 4.0 & 16/16 & 0/16 & 14/16 & 16/16 \\
    No integrity & 5.0 & 5.0 & 0.0 & 5.0 & 4.0 & 16/16 & 16/16 & 0/16 & 16/16 \\
    No receipt & 4.0 & 4.0 & 4.81 & 2.0 & 4.0 & 16/16 & 16/16 & 11/16 & 0/16 \\
    Evidence chain & 5.0 & 5.0 & 4.81 & 5.0 & 5.0 & 16/16 & 16/16 & 11/16 & 16/16 \\
\\
\bottomrule
\end{tabular}%
}
\end{table}


Removing intent capture lowers explicitness from 5.0 to 3.0 and evidence-stage exposure from 5.0 to 4.0, while leaving policy, integrity, and receipt behaviors mostly intact. Removing governance lowers explicitness to 2.0, eliminates all positive policy-checked reviews, and changes the task suite behavior by allowing formerly blocked tasks to complete, while tamper sensitivity remains high because integrity checks are still present. Removing integrity leaves explicitness and audit boundedness at evidence-chain levels, but drives tamper sensitivity and execution verification to 0.0 and 0 of 16. Removing receipt export preserves most other capabilities but drops audit boundedness from 5.0 to 2.0 and receipt-exported reviews from 16 of 16 to 0 of 16.

Taken together, the ablations indicate that each stage contributes a distinct review function. The evidence chain therefore behaves as a layered architecture rather than as a single monolithic switch.

\subsection{Human-review scaffold}

The blinded human-review kit is retained as workflow infrastructure rather than as a main empirical result. Because the checked-in sheet is synthetic, the paper does not use those values as evidence of usability or reviewer performance. We therefore treat the human-review component as protocol readiness only and defer the smoke-test details to Appendix~\ref{sec:appendix-human-review}.
